\documentclass{report}

% --- Core math ---
\usepackage{amsmath, amsfonts, amssymb, mathtools}
\usepackage[varbb]{newpxmath}   % optional: nicer math font; drop if you don't want it
\usepackage{microtype}
\usepackage{xparse}
\usepackage{enumitem}


% --- Layout ---
\usepackage[margin=1in]{geometry}
\usepackage[Glenn]{fncychap}

% --- Colors, boxes, code, references ---
\usepackage{xcolor}
\usepackage{tikz}
\usepackage{varwidth}
\usepackage{listings}
\usepackage[most, many, breakable]{tcolorbox}
\usepackage{hyperref}
\usepackage[nameinlink, noabbrev]{cleveref}

% --- Chemistry ---
\usepackage[version=4]{mhchem}
\usepackage{chemfig}

\hypersetup{
	colorlinks=true,
	linkcolor=doc!90,
}


\hbadness=99999
\title{Lecture Notes}
\author{Mohabeer Taij Ungnoo}
\date{}

%%%%%%%%%%%%%%%%%%%%%%%%%%%%%%%%%%%%%%%%%%%%%%%%%%%%%%%%%%%%%%%%%%%%%%%%%%%%%%
% Much of this code was taken from Charlie's maths template, available here: %
% https://github.com/SeniorMars/dotfiles/tree/main/latex_template            %
%%%%%%%%%%%%%%%%%%%%%%%%%%%%%%%%%%%%%%%%%%%%%%%%%%%%%%%%%%%%%%%%%%%%%%%%%%%%%%


%================================
% COLORS
%================================

\definecolor{doc}{RGB}{0,60,110}
\definecolor{myg}{RGB}{56, 140, 70}        % claim accent
\definecolor{myp}{RGB}{197, 92, 212}       % corollary accent
\definecolor{myb}{RGB}{45, 111, 177}       % question accent

\definecolor{mytheorembg}{HTML}{F2F2F9}
\definecolor{mytheoremfr}{HTML}{00007B}
\definecolor{mylemmabg}{HTML}{FFFAF8}
\definecolor{mylemmafr}{HTML}{983b0f}
\definecolor{myexamplebg}{HTML}{F2FBF8}
\definecolor{myexamplefr}{HTML}{88D6D1}
\definecolor{myexampleti}{HTML}{2A7F7F}
\definecolor{mydefinitbg}{HTML}{FFF5F5}
\definecolor{mydefinitfr}{HTML}{8A1C1C}
\definecolor{myremarkbg}{HTML}{FAF9F9}
\definecolor{myremarkfr}{HTML}{000000}

\definecolor{codebg}{HTML}{F7F8FA}
\definecolor{codegutter}{HTML}{EDEFF3}
\definecolor{codeframe}{HTML}{D8DEE8}
\definecolor{codeaccent}{HTML}{3867D6}
\definecolor{codekeyword}{HTML}{7C3AED}
\definecolor{codestring}{HTML}{0F766E}
\definecolor{codecomment}{HTML}{64748B}
\definecolor{codenumber}{HTML}{94A3B8}
\definecolor{codeidentifier}{HTML}{1F2937}

%================================
% TCOLORBOX BASE STYLES
%================================

\tcbuselibrary{theorems, skins, hooks, listings}

\tcbset{
	cmtleftpanel/.style 2 args={
		enhanced, breakable,
		colback = #1,
		frame hidden,
		boxrule = 0sp,
		borderline west = {2pt}{0pt}{#2},
		sharp corners,
		description font = \mdseries,
		separator sign none,
	},
	cmtleftthm/.style 2 args={
		cmtleftpanel={#1}{#2},
		detach title,
		before upper = \tcbtitle\par\smallskip,
		coltitle = #2,
		fonttitle = \bfseries\sffamily,
		segmentation style={solid, #2},
	},
	cmtexamplethm/.style={
		colback = myexamplebg,
		breakable,
		colframe = myexamplefr,
		coltitle = myexampleti,
		boxrule = 1pt,
		sharp corners,
		detach title,
		before upper=\tcbtitle\par\smallskip,
		fonttitle = \bfseries,
		description font = \mdseries,
		separator sign none,
		description delimiters parenthesis,
	},
	cmtdefinitionthm/.style={
		enhanced,
		before skip=2mm,
		after skip=2mm,
		colback=mydefinitbg,
		colframe=mydefinitfr,
		boxrule=0.5mm,
		attach boxed title to top left={xshift=1cm,yshift*=1mm-\tcboxedtitleheight},
		varwidth boxed title*=-3cm,
		colbacktitle=mydefinitfr,
		coltitle=white,
		boxed title style={
			frame code={
				\path[fill=tcbcolback]
				([yshift=-1mm,xshift=-1mm]frame.north west)
				arc[start angle=0,end angle=180,radius=1mm]
				([yshift=-1mm,xshift=1mm]frame.north east)
				arc[start angle=180,end angle=0,radius=1mm];
				\path[left color=tcbcolback!60!black,right color=tcbcolback!60!black,
					middle color=tcbcolback!80!black]
				([xshift=-2mm]frame.north west) -- ([xshift=2mm]frame.north east)
				[rounded corners=1mm]-- ([xshift=1mm,yshift=-1mm]frame.north east)
				-- (frame.south east) -- (frame.south west)
				-- ([xshift=-1mm,yshift=-1mm]frame.north west)
				[sharp corners]-- cycle;
			},
			interior engine=empty,
		},
		fonttitle=\bfseries,
	},
}

\makeatletter
\tcbset{
	cmtribbonbox/.style={
		enhanced,
		breakable,
		colback=white,
		attach boxed title to top left={yshift*=-\tcboxedtitleheight},
		fonttitle=\bfseries,
		boxed title size=title,
		boxed title style={
			sharp corners,
			rounded corners=northwest,
			colback=tcbcolframe,
			boxrule=0pt,
		},
		underlay boxed title={
			\path[fill=tcbcolframe] (title.south west)--(title.south east)
			to[out=0, in=180] ([xshift=5mm]title.east)--
			(title.center-|frame.east)
			[rounded corners=\kvtcb@arc] |-
			(frame.north) -| cycle;
		},
	},
}
\makeatother

%================================
% THEOREM-FAMILY BOXES
%================================

\newtcbtheorem[number within=section,
	crefname={theorem}{theorems}, Crefname={Theorem}{Theorems}]
	{theorem}{Theorem}{cmtleftthm={mytheorembg}{mytheoremfr}}{th}

\newtcbtheorem[number within=section,
	crefname={corollary}{corollaries}, Crefname={Corollary}{Corollaries}]
	{corollary}{Corollary}{cmtleftthm={myp!10}{myp!85!black}}{th}

\newtcbtheorem[number within=section,
	crefname={lemma}{lemmas}, Crefname={Lemma}{Lemmas}]
	{lemma}{Lemma}{cmtleftthm={mylemmabg}{mylemmafr}}{th}

\newtcbtheorem[number within=section,
	crefname={claim}{claims}, Crefname={Claim}{Claims}]
	{claim}{Claim}{cmtleftthm={myg!10}{myg!85!black}, borderline west={2pt}{0pt}{myg}}{th}

\newtcbtheorem[number within=section,
	crefname={example}{examples}, Crefname={Example}{Examples}]
	{example}{Example}{cmtexamplethm}{ex}

\newtcbtheorem[number within=section,
	crefname={definition}{definitions}, Crefname={Definition}{Definitions}]
	{definition}{Definition}{cmtdefinitionthm, title={#2}, #1}{def}

\newtcbtheorem[number within=section,
	crefname={remark}{remarks}, Crefname={Remark}{Remarks}]
	{remark}{Remark}{cmtleftthm={myremarkbg}{myremarkfr}}{rmk}

\newtcbtheorem[number within=section,
	crefname={question}{questions}, Crefname={Question}{Questions}]
	{question}{Question}{cmtribbonbox, colframe=myb!80!black, title={#2}, #1}{qst}

%================================
% NOTE BOX (unnumbered)
%================================

\usetikzlibrary{arrows, calc, shadows.blur}

\newtcolorbox{quotebox}[1][]{
	cmtleftpanel={myp!10}{myp!85!black},
	fontupper=\itshape,
	#1
}
\newtcolorbox{note}[1][]{
	enhanced jigsaw,
	colback=gray!20!white,
	colframe=gray!80!black,
	size=small,
	boxrule=1pt,
	title=\textbf{Note:-},
	halign title=flush center,
	coltitle=black,
	breakable,
	drop shadow=black!50!white,
	attach boxed title to top left={xshift=1cm,yshift=-\tcboxedtitleheight/2,yshifttext=-\tcboxedtitleheight/2},
	minipage boxed title=1.5cm,
	boxed title style={
		colback=white,
		size=fbox,
		boxrule=1pt,
		boxsep=2pt,
		underlay={
			\coordinate (dotA) at ($(interior.west) + (-0.5pt,0)$);
			\coordinate (dotB) at ($(interior.east) + (0.5pt,0)$);
			\begin{scope}
				\clip (interior.north west) rectangle ([xshift=3ex]interior.east);
				\filldraw [white, blur shadow={shadow opacity=60, shadow yshift=-.75ex}, rounded corners=2pt]
					(interior.north west) rectangle (interior.south east);
			\end{scope}
			\begin{scope}[gray!80!black]
				\fill (dotA) circle (2pt);
				\fill (dotB) circle (2pt);
			\end{scope}
		},
	},
}

%================================
% CODEBLOCK
%================================

\lstdefinestyle{cmtcodebase}{
	basicstyle=\ttfamily\footnotesize,
	keywordstyle=\bfseries\color{codekeyword},
	commentstyle=\itshape\color{codecomment},
	stringstyle=\color{codestring},
	identifierstyle=\color{codeidentifier},
	numberstyle=\scriptsize\color{codenumber},
	showstringspaces=false,
	tabsize=4,
	columns=fullflexible,
	keepspaces=true,
	upquote=true,
	breaklines=true,
	breakatwhitespace=false,
	postbreak=\mbox{\textcolor{codenumber}{$\hookrightarrow$}\space},
}
\lstdefinestyle{cmtcode}{
	style=cmtcodebase,
	numbers=left,
	numbersep=10pt,
	xleftmargin=2.8em,
}

\tcbset{
	cmtcodebox/.style={
		enhanced, breakable, sharp corners,
		boxrule=0pt, frame hidden,
		colback=codebg, colframe=codeframe,
		coltitle=codeidentifier,
		fonttitle=\bfseries\sffamily\small,
		colbacktitle=codegutter,
		titlerule=0pt,
		toptitle=4pt, bottomtitle=4pt,
		lefttitle=8pt, righttitle=8pt,
		left=8pt, right=8pt, top=6pt, bottom=6pt,
		boxsep=0pt,
		before skip=\baselineskip, after skip=\baselineskip,
		borderline west={2pt}{0pt}{codeaccent},
		segmentation hidden,
		listing only,
	},
}
\newtcblisting{codeblock}[2][]{
	cmtcodebox,
	title={#2},
	listing options={style=cmtcode, language=#2},
	#1
}

%================================
% SHORT COMMANDS
%================================

\newcommand{\dfn}[2]{\begin{definition}{#1}{}#2\end{definition}}
\newcommand{\ex}[2]{\begin{example}{#1}{}#2\end{example}}
\newcommand{\rmk}[2]{\begin{remark}{#1}{}#2\end{remark}}
\newcommand{\clm}[3]{\begin{claim}{#1}{#2}#3\end{claim}}
\newcommand{\thm}[2]{\begin{theorem}{#1}{}#2\end{theorem}}
\newcommand{\cor}[2]{\begin{corollary}{#1}{}#2\end{corollary}}
\newcommand{\mlemma}[2]{\begin{lemma}{#1}{}#2\end{lemma}}
\newcommand{\qst}[2]{\begin{question}{#1}{}#2\end{question}}
\newcommand{\nt}[1]{\begin{note}#1\end{note}}
\newcommand{\qte}[1]{\begin{quotebox}#1\end{quotebox}}


% Ordinary derivative: \dv[n]{x}{f} -> d^n f / dx^n ; \dv[n]{x} -> bare operator d^n/dx^n
\NewDocumentCommand{\dv}{o m g}{%
	\IfValueTF{#3}
		{\IfNoValueTF{#1}
			{\frac{\mathrm{d} #3}{\mathrm{d} #2}}
			{\frac{\mathrm{d}^{#1} #3}{\mathrm{d} #2^{#1}}}}
		{\IfNoValueTF{#1}
			{\frac{\mathrm{d}}{\mathrm{d} #2}}
			{\frac{\mathrm{d}^{#1}}{\mathrm{d} #2^{#1}}}}
}

% Partial derivative: \ddv[n]{x}{f} -> partial^n f / partial x^n ; \ddv[n]{x} -> bare operator
\NewDocumentCommand{\ddv}{o m g}{%
	\IfValueTF{#3}
		{\IfNoValueTF{#1}
			{\frac{\partial #3}{\partial #2}}
			{\frac{\partial^{#1} #3}{\partial #2^{#1}}}}
		{\IfNoValueTF{#1}
			{\frac{\partial}{\partial #2}}
			{\frac{\partial^{#1}}{\partial #2^{#1}}}}
}

\newcommand{\m}{\mathrm{m}}
\newcommand{\mol}{\mathrm{mol}}
\newcommand{\cm}{\mathrm{cm}}
\newcommand{\kg}{\mathrm{kg}}
\newcommand{\K}{\mathrm{K}}
\newcommand{\s}{\mathrm{s}}


% Not yet in template.tex - added here for this reference doc
\usepackage[version=4]{mhchem}
\usepackage{chemfig}

\begin{document}
\maketitle

\section{Theorem-family boxes}

\thm{Fundamental Theorem of Calculus}{
}

\cor{Constant Antiderivatives}{
	If $F$ and $G$ are both antiderivatives of $f$ on an interval $I$,
	then $F - G$ is constant on $I$.
}

\mlemma{Bounded Monotone Sequences}{
	Every bounded, monotone sequence of real numbers converges.
}

\clm{Sum of Two Evens}{is even}{
	Let $a, b$ be even integers. Then $a + b$ is even, since
	$a = 2m$ and $b = 2n$ for some $m, n \in \mathbb{Z}$, so
	$a + b = 2(m+n)$.
}

\ex{Evaluating a Limit}{
	Compute $\displaystyle\lim_{x \to 0} \frac{\sin x}{x}$.
	By the squeeze theorem, this limit equals $1$.
}

\dfn{Continuity}{
	A function $f$ is \emph{continuous} at $a$ if
	$\displaystyle\lim_{x \to a} f(x) = f(a)$.
}

\rmk{On Notation}{
	See above for the formal definition of continuity used
	throughout these notes.
}

\qst{Open Problem}{
	Does the Collatz conjecture hold for all positive integers?
}

\nt{
	Boxes that share the \texttt{cmtleftthm} style (theorem, corollary, lemma,
	claim, remark) are numbered within the section and share the \texttt{th}
	label family; \texttt{definition} and \texttt{question} use their own
	\texttt{def}/\texttt{qst} prefixes instead.
}

\begin{codeblock}{Python}
def is_prime(n):
    if n < 2:
        return False
    for k in range(2, int(n**0.5) + 1):
        if n % k == 0:
            return False
    return True
\end{codeblock}

\section{Derivative macros}

Ordinary derivative, with and without an explicit function:
\begin{itemize}
	\item \verb|\dv{x}{f}| \quad $\to$ \quad $\dv{x}{f}$
	\item \verb|\dv{x}| \quad $\to$ \quad $\dv{x}\left(x^2 + 3x\right)$
	\item \verb|\dv[2]{x}{f}| \quad $\to$ \quad $\dv[2]{x}{f}$
	\item \verb|\dv[3]{x}| \quad $\to$ \quad $\dv[3]{x}\left(\sin x\right)$
\end{itemize}

Partial derivative, same four forms:
\begin{itemize}
	\item \verb|\ddv{x}{f}| \quad $\to$ \quad $\ddv{x}{f}$
	\item \verb|\ddv{y}| \quad $\to$ \quad $\ddv{y}\left(x^2 y + \sin(y)\right)$
	\item \verb|\ddv[2]{x}{\sin(x) + 7y}| \quad $\to$ \quad $\ddv[2]{x}{\sin(x) + 7y}$
	\item \verb|\ddv[2]{y}| \quad $\to$ \quad $\ddv[2]{y}\left(x^2 y^3\right)$
\end{itemize}

A mixed use, e.g. writing out a Laplacian in two variables:
\[
	\ddv[2]{x}{u} + \ddv[2]{y}{u} = 0
\]

\section{mhchem}

\subsection{Basic formulas and reactions}

Simple molecular formulas: \ce{H2O}, \ce{CO2}, \ce{C6H12O6}.

Balanced reaction with an arrow:
\[
	\ce{2H2 + O2 -> 2H2O}
\]

Reversible reaction (equilibrium arrows):
\[
	\ce{CO2 + H2O <=> H2CO3}
\]

Reaction with states of matter:
\[
	\ce{Na(s) + H2O(l) -> NaOH(aq) + 1/2 H2(g)}
\]

\subsection{Charges and isotopes}

Ions: \ce{Na+}, \ce{Cl-}, \ce{SO4^2-}, \ce{Fe^3+}.

Precipitation reaction:
\[
	\ce{Ag+ + Cl- -> AgCl v}
\]

Isotope notation:
\[
	\ce{^{14}_6C ->[\beta^-] \, ^{14}_7N}
\]

\subsection{Equilibrium constant}

\[
	\ce{N2 + 3H2 <=>[$K_c$] 2NH3}
\]

\section{chemfig}

\subsection{Simple chains}

Ethanol, drawn as a chain of bonded groups:
\[
	\chemfig{CH_3-CH_2-OH}
\]

Propane:
\[
	\chemfig{CH_3-CH_2-CH_3}
\]

\subsection{Double and triple bonds}

Ethene (double bond):
\[
	\chemfig{H_2C=CH_2}
\]

Ethyne (triple bond):
\[
	\chemfig{HC~CH}
\]

\subsection{Rings}

Benzene, drawn as a hexagon with alternating double bonds:
\[
	\chemfig{*6(=-=-=-)}
\]

Benzene, drawn with the circle:
\[
	\chemfig{**6(------)}
\]

\subsection{Branching}

Isobutane, using parentheses to branch off the main chain:
\[
	\chemfig{CH_3-CH(-[2]CH_3)-CH_3}
\]

\begin{figure}
\centering
\begin{tikzpicture}[scale=1.1]
  % Axes
  \draw[->] (-0.2,-3.2) -- (-0.2,3.0) node[above] {$PE$};
  \draw[->] (-0.2,0) -- (7.2,0);
  \node[right] at (5.6,0.6) {\textit{Distance of}};
  \node[right] at (5.6,0.1) {\textit{separation}};

   \latticelabels{0}{6}{-3}{3}

  % Tick labels
  \node[left] at (-0.3,0) {$0$};
  \node[left] at (-0.3,1.8) {$+$};
  \node[left] at (-0.3,-1.8) {$-$};

  % Curve defined by points and tangent angles
  % P1 (0.0,2.8): leaves steeply downward (-80 deg)
  % P2 (1.6,-2.5): horizontal tangent (minimum), arrives at 180 and leaves at 0
  % P3 (6.6,-0.05): arrives horizontally from the left
  \draw[thick] (0.0,2.8)
    to[out=-80, in=180] (1.6,-2.5)
    to[out=0,   in=180] (6.6,-0.05);
\end{tikzpicture}
\caption{Potential energy as a function of distance of separation between two molecules.}
\label{fig:pe}
\end{figure}

\begin{figure}
\centering
\begin{tikzpicture}
  \begin{axis}[
    axis lines=middle,
    xlabel={$x$}, ylabel={$\sin x$},
    domain=0:2*pi, samples=100,
    xtick={0,pi/2,pi,3*pi/2,2*pi},
    xticklabels={$0$,$\frac{\pi}{2}$,$\pi$,$\frac{3\pi}{2}$,$2\pi$},
    ymin=-1.3, ymax=1.3 ]
    \addplot[thick, blue] {sin(deg(x))};
  \end{axis}
\end{tikzpicture}
\caption{The sine function.}
\end{figure}

\begin{figure}[ht]
\centering
\begin{tikzpicture}
  \begin{axis}[
    axis lines=middle,
    xlabel={$V$}, ylabel={$p$},
    domain=0:2, samples=100,
    ymin=0, ymax=4,
    xmin=0, xmax = 4,
    xtick = 1.5, ytick = 2,
    xticklabels = {$V_C$},
    yticklabels = {$P_C$},
    clip = false
    ]
    \draw[name path=pres,dashed]    (.4,.2) .. controls ++(.2,1) and ++(-.5,0) ..
      (1.5,2) node[below]{X} node[circle, fill, inner sep = 1.5]{} .. controls ++(.5,0) and ++(-1.5,.5) .. 
      (3.8,.2) ;
    \path[name path = L] (0,1.2) --++ (4,0);
    \path[name intersections = {of = L and pres}];
    \coordinate (B) at (intersection-1);
    \coordinate (A) at (intersection-2);
    \draw (B) node[circle, fill, inner sep = 1.5pt]{} node[left]{B}-- (A) node[circle, fill, inner sep = 1.5pt]{} node[right]{A};
    \draw[dashed] (1.5,0) -- (1.5, 2);
    \draw[dashed] (0,2) -- (1.5, 2);

    \draw[blue] (0.5, 3.5) .. controls ++ (0.1, -1.5) and ++(-0.2, 0) .. (1.5, 2) .. controls ++(0.5, 0) and ++(-1, 0.2) .. (3.8, 0.7) node at ++(0:1em){$T_C$};
    \draw (0.4, 3.5) to[bend right = 7] (B) -- (A) to[bend right = 10] (3.8, 0.4) node[right]{$T_1$};
    \draw (0.6, 3.5) to[bend right = 17] (3.8, 1) node[right]{$T_2$};
  \end{axis}
\end{tikzpicture}
\caption{Multiple isotherms of a typical real gas.}
\end{figure}

\end{document}
